\documentclass[aspectratio=169,11pt]{beamer}

\usetheme{AnnArbor}
\usecolortheme{dolphin}

\usepackage{fontspec}
\usepackage{polyglossia}
\setdefaultlanguage{english}
\setmainfont{DejaVu Serif}
\setsansfont{DejaVu Sans}
\setmonofont{DejaVu Sans Mono}

\usepackage{booktabs}
\usepackage{tabularx}
\usepackage{array}
\usepackage{xcolor}
\usepackage{listings}
\usepackage{tikz}
\usetikzlibrary{arrows.meta,positioning}

\definecolor{sqlbg}{RGB}{248,248,248}
\definecolor{sqlkw}{RGB}{0,70,140}
\definecolor{sqlcm}{RGB}{90,90,90}
\definecolor{myblue}{RGB}{32,83,150}
\definecolor{myorange}{RGB}{202,100,30}

\lstdefinestyle{sqlstyle}{
  language=SQL,
  basicstyle=\ttfamily\scriptsize,
  keywordstyle=\color{sqlkw}\bfseries,
  commentstyle=\color{sqlcm}\itshape,
  stringstyle=\color{myorange},
  backgroundcolor=\color{sqlbg},
  frame=single,
  rulecolor=\color{black!20},
  columns=fullflexible,
  keepspaces=true,
  showstringspaces=false,
  breaklines=true,
  tabsize=2
}
\lstset{style=sqlstyle}

\newcommand{\code}[1]{\texttt{#1}}
\newcommand{\smallitem}{\setlength{\itemsep}{2pt}}
\newcommand{\tinyitem}{\setlength{\itemsep}{1pt}}

\title[MySQL Storage Engines]{MySQL Storage Engines}
\subtitle{Concepts, Verification, Selection, and Practical Applications}
\author{DSAI1002 -- Database Tutorial}
\date{Updated: June 4, 2026}

\begin{document}

% 1
\begin{frame}
  \titlepage
\end{frame}

% 2
\begin{frame}{Learning Objectives}
\begin{columns}[T,onlytextwidth]
\column{0.48\textwidth}
\begin{block}{After this lesson, learners can}
\begin{enumerate}\smallitem
  \item Explain the concept of a \alert{storage engine} in MySQL.
  \item Describe the engine's role in storing and processing table data.
  \item Inspect available engines using \code{SHOW ENGINES} and \code{information\_schema.engines}.
\end{enumerate}
\end{block}
\column{0.48\textwidth}
\begin{block}{Practical Skills}
\begin{enumerate}\smallitem
  \setcounter{enumi}{3}
  \item Specify an engine when creating tables using \code{ENGINE}.
  \item Differentiate commonly used storage engines.
  \item Choose the appropriate engine for real-world scenarios.
\end{enumerate}
\end{block}
\end{columns}
\end{frame}

% 3
\begin{frame}{High-Level Overview}
\begin{columns}[T,onlytextwidth]
\column{0.50\textwidth}
\begin{block}{In MySQL}
Each table is managed by a specific \alert{storage engine}.
\end{block}
\begin{itemize}\smallitem
  \item The engine governs how table data is stored, read, written, locked, and recovered.
  \item Not all engines support the same set of capabilities.
  \item For example: \code{InnoDB} supports transactions and foreign keys; \code{MEMORY} stores data entirely in RAM.
\end{itemize}
\column{0.46\textwidth}
\begin{block}{Two Primary Layers}
\begin{description}\smallitem
  \item[SQL Layer] Receives statements, parses syntax, and optimizes query execution plans.
  \item[Engine Layer] Manages physical storage and data retrieval for each table.
\end{description}
\end{block}
\end{columns}
\end{frame}

% 4
\begin{frame}{MySQL Server and Storage Engines}
\centering
\begin{tikzpicture}[node distance=0.95cm, every node/.style={font=\small}]
\node[draw, rounded corners, fill=blue!8, minimum width=7.9cm, minimum height=0.70cm] (client) {Client / Application sends SQL queries};
\node[draw, rounded corners, fill=green!8, below=of client, minimum width=7.9cm, minimum height=0.70cm] (sql) {SQL Layer: parser, optimizer, execution coordinator};
\node[draw, rounded corners, fill=orange!10, below=of sql, minimum width=7.9cm, minimum height=0.70cm] (engine) {Storage Engine of the target table};
\node[draw, rounded corners, fill=gray!12, below=of engine, minimum width=7.9cm, minimum height=0.70cm] (data) {Physical Data: pages, indexes, logs, files, RAM};
\draw[-{Latex[length=2.4mm]}] (client) -- (sql);
\draw[-{Latex[length=2.4mm]}] (sql) -- (engine);
\draw[-{Latex[length=2.4mm]}] (engine) -- (data);
\end{tikzpicture}
\vspace{0.15cm}

\footnotesize Storage engines operate beneath the SQL layer and directly manage table-level data storage.
\end{frame}

% 5
\begin{frame}{What is a Storage Engine?}
\begin{block}{Definition}
A \alert{storage engine} is a software module within MySQL responsible for handling data storage, retrieval, and management for individual tables.
\end{block}
\begin{columns}[T,onlytextwidth]
\column{0.48\textwidth}
\begin{block}{Determines Storage Format}
\begin{itemize}\smallitem
  \item Whether data resides on disk or in system memory.
  \item How indexing, caching, and crash recovery are organized.
\end{itemize}
\end{block}
\column{0.48\textwidth}
\begin{block}{Determines Feature Support}
\begin{itemize}\smallitem
  \item Whether transactions (ACID) are supported.
  \item Whether foreign keys and row-level locking are available.
\end{itemize}
\end{block}
\end{columns}
\end{frame}

% 6
\begin{frame}[fragile]{Default Storage Engine}
\begin{itemize}\smallitem
  \item Since MySQL 5.5, \code{InnoDB} has been the default storage engine.
  \item If a \code{CREATE TABLE} statement omits the \code{ENGINE} clause, MySQL applies the server default.
\end{itemize}
\begin{lstlisting}
CREATE TABLE students (
    student_id INT PRIMARY KEY,
    student_name VARCHAR(100)
);
\end{lstlisting}
\begin{block}{Interpretation}
If the default engine is \code{InnoDB}, the \code{students} table will be created using \code{InnoDB} automatically.
\end{block}
\end{frame}

% 7
\begin{frame}[fragile]{Specifying the Storage Engine at Creation}
\begin{block}{General Syntax}
\begin{lstlisting}
CREATE TABLE table_name (
    column_list
) ENGINE = engine_name;
\end{lstlisting}
\end{block}
\begin{block}{Example: Creating a Table with InnoDB}
\begin{lstlisting}
CREATE TABLE customers (
    customer_id INT PRIMARY KEY,
    customer_name VARCHAR(100),
    phone VARCHAR(20)
) ENGINE = InnoDB;
\end{lstlisting}
\end{block}
\end{frame}

% 8
\begin{frame}{SQL Processing Flow with Storage Engines}
\begin{enumerate}\smallitem
  \item \textbf{Receive SQL statement:} Client issues \code{SELECT}, \code{INSERT}, \code{UPDATE}, \code{DELETE}, or \code{CREATE TABLE}.
  \item \textbf{Parse and optimize:} MySQL validates syntax, permissions, and formulates an execution plan.
  \item \textbf{Call storage engine:} MySQL delegates read/write requests down to the table's designated engine.
  \item \textbf{Access data:} The engine reads, writes, locks, recovers, or indexes data via its internal mechanisms.
  \item \textbf{Return results:} MySQL receives the data records from the engine and returns the response to the client.
\end{enumerate}
\end{frame}

% 9
\begin{frame}[fragile]{Checking Available Storage Engines}
\begin{columns}[T,onlytextwidth]
\column{0.48\textwidth}
\begin{block}{Method 1: \code{SHOW ENGINES}}
\begin{lstlisting}
SHOW ENGINES;
\end{lstlisting}
\end{block}
\column{0.48\textwidth}
\begin{block}{Method 2: System Table}
\begin{lstlisting}
SELECT
    engine,
    support
FROM information_schema.engines
ORDER BY engine;
\end{lstlisting}
\end{block}
\end{columns}
\vspace{0.25cm}
\begin{alertblock}{Key Takeaway}
These queries identify which engines are supported, which one is set as default, or if any are disabled.
\end{alertblock}
\end{frame}

% 10
\begin{frame}{Support Values in \code{SHOW ENGINES}}
\begin{center}
\begin{tabularx}{0.88\textwidth}{>{\ttfamily}l X}
\toprule
\textbf{Value} & \textbf{Meaning} \\
\midrule
YES & Engine is supported on this server. \\
NO & Engine is not supported. \\
DEFAULT & Engine is supported and currently configured as default. \\
DISABLED & Engine exists in the binary but is currently disabled. \\
\bottomrule
\end{tabularx}
\end{center}
\begin{block}{Key Evaluation Criteria}
Transactions, foreign keys, crash recovery, persistence, and read/write optimization profiles.
\end{block}
\end{frame}

% 11
\begin{frame}{Major Storage Engines in MySQL}
\begin{columns}[T,onlytextwidth]
\column{0.48\textwidth}
\begin{block}{Common Engines}
\begin{itemize}\smallitem
  \item \code{InnoDB}: Default engine, transaction-safe (ACID).
  \item \code{MyISAM}: Legacy engine, read-heavy, non-transactional.
  \item \code{MEMORY}: Fast in-memory volatile storage.
  \item \code{ARCHIVE}: Compressed historical archive storage.
\end{itemize}
\end{block}
\column{0.48\textwidth}
\begin{block}{Specialized Engines}
\begin{itemize}\smallitem
  \item \code{CSV}: Stores data as plain text CSV files.
  \item \code{BLACKHOLE}: Accepts incoming data without storing it.
  \item \code{MERGE/MRG\_MyISAM}: Unions identical MyISAM tables.
  \item \code{FEDERATED}: Connects to remote MySQL servers.
\end{itemize}
\end{block}
\end{columns}
\end{frame}

% 12
\begin{frame}{InnoDB: The Default and Standard Engine}
\begin{block}{Concept}
\code{InnoDB} is the safest and most capable choice for enterprise systems due to its robust transactional support, referential integrity, and crash resilience.
\end{block}
\begin{columns}[T,onlytextwidth]
\column{0.24\textwidth}
\begin{block}{Transactions}
\small Supports \code{COMMIT}, \code{ROLLBACK}, and ACID properties.
\end{block}
\column{0.24\textwidth}
\begin{block}{Foreign Keys}
\small Enforces referential integrity constraints.
\end{block}
\column{0.24\textwidth}
\begin{block}{Row Locking}
\small Reduces lock contention in concurrent multi-user workloads.
\end{block}
\column{0.24\textwidth}
\begin{block}{Crash Recovery}
\small Automatic recovery using write-ahead redo logging.
\end{block}
\end{columns}
\end{frame}

% 13
\begin{frame}[fragile]{Example: Creating an \code{InnoDB} Table}
\begin{lstlisting}
CREATE TABLE orders (
    order_id INT PRIMARY KEY,
    customer_id INT NOT NULL,
    order_date DATE,
    total_amount DECIMAL(12, 2)
) ENGINE = InnoDB;
\end{lstlisting}
\begin{block}{Recommended Use Cases for \code{InnoDB}}
\begin{itemize}\smallitem
  \item E-commerce, billing, and sales order processing systems.
  \item Financial accounting, payments, and user identity tables.
  \item Any table demanding strict data integrity, foreign keys, and error rollbacks.
\end{itemize}
\end{block}
\end{frame}

% 14
\begin{frame}[fragile]{MyISAM}
\begin{columns}[T,onlytextwidth]
\column{0.50\textwidth}
\begin{block}{Characteristics}
\begin{itemize}\smallitem
  \item Default engine prior to MySQL 5.5.
  \item Can be fast in specific read-heavy scenarios.
  \item Does not support ACID transactions.
  \item Does not support foreign key constraints.
\end{itemize}
\end{block}
\column{0.46\textwidth}
\begin{block}{Example}
\begin{lstlisting}
CREATE TABLE article_archive (
    article_id INT PRIMARY KEY,
    title VARCHAR(255),
    content TEXT
) ENGINE = MyISAM;
\end{lstlisting}
\end{block}
\end{columns}
\end{frame}

% 15
\begin{frame}[fragile]{MEMORY}
\begin{columns}[T,onlytextwidth]
\column{0.50\textwidth}
\begin{block}{Characteristics}
\begin{itemize}\smallitem
  \item Stores all table data in RAM.
  \item Ideal for temporary caching and lookup tables.
  \item Extremely low-latency reads and writes.
  \item Volatile: data is wiped when MySQL Server stops or restarts.
\end{itemize}
\end{block}
\column{0.46\textwidth}
\begin{block}{Example}
\begin{lstlisting}
CREATE TABLE session_cache (
    session_id VARCHAR(100) PRIMARY KEY,
    user_id INT,
    created_at DATETIME
) ENGINE = MEMORY;
\end{lstlisting}
\end{block}
\end{columns}
\end{frame}

% 16
\begin{frame}[fragile]{ARCHIVE and CSV}
\begin{columns}[T,onlytextwidth]
\column{0.48\textwidth}
\begin{block}{ARCHIVE}
\small Stores large volumes of unindexed, append-only historical records in a compressed format; ideal for audit logs.
\begin{lstlisting}
CREATE TABLE audit_logs (
    log_id INT,
    action_name VARCHAR(100),
    action_time DATETIME
) ENGINE = ARCHIVE;
\end{lstlisting}
\end{block}
\column{0.48\textwidth}
\begin{block}{CSV}
\small Stores data directly in text files using comma-separated values; easy for file exchange but lacks indexing and constraints.
\begin{lstlisting}
CREATE TABLE export_customers (
    customer_id INT NOT NULL,
    customer_name VARCHAR(100) NOT NULL,
    phone VARCHAR(20) NOT NULL
) ENGINE = CSV;
\end{lstlisting}
\end{block}
\end{columns}
\end{frame}

% 17
\begin{frame}{BLACKHOLE, MERGE, and FEDERATED}
\begin{columns}[T,onlytextwidth]
\column{0.32\textwidth}
\begin{block}{BLACKHOLE}
\small Discards written data while maintaining binary logs. Used in replication relaying and testing.
\end{block}
\column{0.32\textwidth}
\begin{block}{MERGE / MRG\_MyISAM}
\small Combines multiple identically structured \code{MyISAM} tables into one logical view.
\end{block}
\column{0.32\textwidth}
\begin{block}{FEDERATED}
\small Queries tables on remote MySQL instances without holding physical table data locally.
\end{block}
\end{columns}
\vspace{0.25cm}
\begin{alertblock}{Note}
These engines serve specific architectural needs and are not intended for general application tables.
\end{alertblock}
\end{frame}

% 18
\begin{frame}{Principles for Choosing a Storage Engine}
\begin{columns}[T,onlytextwidth]
\column{0.48\textwidth}
\begin{block}{Data Integrity Requirements}
If tables have foreign key relations, require rollbacks, or process financial data, select \code{InnoDB}.
\end{block}
\begin{block}{Data Durability Needs}
If records must persist across server restarts, never use \code{MEMORY} for persistent data.
\end{block}
\column{0.48\textwidth}
\begin{block}{Performance Goals}
Temporary calculation tables benefit from \code{MEMORY}; robust concurrent read/write workloads belong on \code{InnoDB}.
\end{block}
\begin{block}{Nature of Data}
High-volume write-once audit logs suit \code{ARCHIVE}; external exports may use \code{CSV}.
\end{block}
\end{columns}
\end{frame}

% 19
\begin{frame}{Practical Use Cases}
\begin{center}
\begin{tabularx}{0.96\textwidth}{p{0.28\textwidth} X}
\toprule
\textbf{Scenario} & \textbf{Suitable Engine} \\
\midrule
Sales \& E-commerce & \code{InnoDB}: Demands transactions, rollback, foreign keys, and ACID guarantees. \\
Temporary Caching & \code{MEMORY}: Fast short-lived intermediate data storage. \\
Log Archiving & \code{ARCHIVE}: High-volume historical records, minimal updates, compact storage. \\
Data Interchange & \code{CSV}: Easily read by spreadsheet applications and external export scripts. \\
Special Testing & \code{BLACKHOLE}: Captures write commands without consuming disk capacity. \\
\bottomrule
\end{tabularx}
\end{center}
\end{frame}

% 20
\begin{frame}{Role in Database Administration}
\begin{columns}[T,onlytextwidth]
\column{0.48\textwidth}
\begin{block}{Database Design}
Select engines tailored to each table's role; primary business tables should consistently use \code{InnoDB}.
\end{block}
\begin{block}{Performance Tuning}
Understanding engine mechanics assists in diagnosing slow queries, lock contention, and I/O bottlenecks.
\end{block}
\column{0.48\textwidth}
\begin{block}{Backup and Recovery}
Different engines use distinct on-disk formats, requiring customized backup and restore workflows.
\end{block}
\begin{block}{System Operations}
During version upgrades or server migrations, DBAs must verify supported and default engines.
\end{block}
\end{columns}
\end{frame}

% 21
\begin{frame}{Quick Comparison: Common Engines}
\scriptsize
\begin{tabularx}{0.95\textwidth}{p{0.24\textwidth}*{4}{>{\centering\arraybackslash}X}}
\toprule
\textbf{Criteria} & \textbf{InnoDB} & \textbf{MyISAM} & \textbf{MEMORY} & \textbf{ARCHIVE} \\
\midrule
Transactions & Yes & No & No & No \\
Foreign Keys & Yes & No & No & No \\
Data Durability & Yes & Yes & No & Yes \\
Primary Business Data & Yes & Limited & No & Limited \\
Primary Use Case & OLTP / Business & Read-heavy / legacy & Staging tables & Logs / history \\
\bottomrule
\end{tabularx}
\vspace{0.35cm}
\begin{block}{Practical Guideline}
For the vast majority of modern applications, \code{InnoDB} is the safest and recommended default.
\end{block}
\end{frame}

% 22
\begin{frame}{Quick Comparison: Specialized Engines}
\scriptsize
\begin{tabularx}{0.95\textwidth}{p{0.27\textwidth}*{3}{>{\centering\arraybackslash}X}}
\toprule
\textbf{Criteria} & \textbf{CSV} & \textbf{BLACKHOLE} & \textbf{FEDERATED} \\
\midrule
Transactions & No & No & No \\
Foreign Keys & No & No & No \\
Data Durability & Yes & No & Depends on remote server \\
Primary Business Data & No & No & Limited \\
Primary Use Case & CSV file export & Discard writes / replication & Remote table access \\
\bottomrule
\end{tabularx}
\vspace{0.35cm}
\begin{alertblock}{Caution}
Specialized engines should only be deployed when there is an explicit architectural justification.
\end{alertblock}
\end{frame}

% 22
\begin{frame}{Quick Quiz 1}
\begin{block}{Question 1}
To which entity is a MySQL storage engine directly attached?
\end{block}
\vspace{0.2cm}
\begin{columns}[T,onlytextwidth]
\column{0.48\textwidth}
\begin{itemize}\smallitem
  \item[A.] User account
  \item[B.] Table
\end{itemize}
\column{0.48\textwidth}
\begin{itemize}\smallitem
  \item[C.] Database password
  \item[D.] SQL comment
\end{itemize}
\end{columns}
\end{frame}

% 23
\begin{frame}{Quick Quiz 2}
\begin{block}{Question 2}
Why is understanding storage engines essential?
\end{block}
\vspace{0.2cm}
\begin{itemize}\smallitem
  \item[A.] Because the engine defines the entire SQL syntax.
  \item[B.] Because each engine has distinct storage characteristics and capabilities.
  \item[C.] Because MySQL supports only a single storage engine.
  \item[D.] Because engines are only used for renaming tables.
\end{itemize}
\end{frame}

% 24
\begin{frame}{Quick Quiz 3}
\begin{block}{Question 3}
Which storage engine is default in modern MySQL releases?
\end{block}
\vspace{0.2cm}
\begin{columns}[T,onlytextwidth]
\column{0.48\textwidth}
\begin{itemize}\smallitem
  \item[A.] MyISAM
  \item[B.] MEMORY
\end{itemize}
\column{0.48\textwidth}
\begin{itemize}\smallitem
  \item[C.] InnoDB
  \item[D.] CSV
\end{itemize}
\end{columns}
\end{frame}

% 25
\begin{frame}{Quick Quiz 4}
\begin{block}{Question 4}
Which clause is used to specify the storage engine when creating a table?
\end{block}
\vspace{0.2cm}
\begin{columns}[T,onlytextwidth]
\column{0.48\textwidth}
\begin{itemize}\smallitem
  \item[A.] \code{USING}
  \item[B.] \code{ENGINE}
\end{itemize}
\column{0.48\textwidth}
\begin{itemize}\smallitem
  \item[C.] \code{STORAGE PATH}
  \item[D.] \code{TABLE TYPE ONLY}
\end{itemize}
\end{columns}
\end{frame}

% 26
\begin{frame}{Quick Quiz 5}
\begin{block}{Question 5}
Which command lists all available storage engines on the server?
\end{block}
\vspace{0.2cm}
\begin{columns}[T,onlytextwidth]
\column{0.48\textwidth}
\begin{itemize}\smallitem
  \item[A.] \code{SHOW TABLES;}
  \item[B.] \code{SHOW ENGINES;}
\end{itemize}
\column{0.48\textwidth}
\begin{itemize}\smallitem
  \item[C.] \code{SHOW DATABASES;}
  \item[D.] \code{SHOW USERS;}
\end{itemize}
\end{columns}
\end{frame}

% 27
\begin{frame}{Quick Quiz 6}
\begin{block}{Question 6}
What does the value \code{DEFAULT} signify in \code{SHOW ENGINES}?
\end{block}
\vspace{0.2cm}
\begin{itemize}\smallitem
  \item[A.] The engine is disabled.
  \item[B.] The engine does not exist.
  \item[C.] The engine is currently configured as the server default.
  \item[D.] The engine is used exclusively for backups.
\end{itemize}
\end{frame}

% 28
\begin{frame}{Quick Quiz 7}
\begin{block}{Question 7}
Which type of system is \code{InnoDB} best suited for?
\end{block}
\vspace{0.2cm}
\begin{itemize}\smallitem
  \item[A.] Systems requiring transactions, rollbacks, and foreign keys.
  \item[B.] Tables used only to discard write data.
  \item[C.] CSV files opened directly in Excel.
  \item[D.] Ephemeral tables that lose data upon server shutdown.
\end{itemize}
\end{frame}

% 29
\begin{frame}{Quick Quiz 8}
\begin{block}{Question 8}
Which storage engine stores all its data directly in RAM?
\end{block}
\vspace{0.2cm}
\begin{columns}[T,onlytextwidth]
\column{0.48\textwidth}
\begin{itemize}\smallitem
  \item[A.] MEMORY
  \item[B.] CSV
\end{itemize}
\column{0.48\textwidth}
\begin{itemize}\smallitem
  \item[C.] BLACKHOLE
  \item[D.] FEDERATED
\end{itemize}
\end{columns}
\end{frame}

% 30
\begin{frame}{Quick Quiz 9}
\begin{block}{Question 9}
Which engine receives data without persisting it to disk or memory?
\end{block}
\vspace{0.2cm}
\begin{columns}[T,onlytextwidth]
\column{0.48\textwidth}
\begin{itemize}\smallitem
  \item[A.] InnoDB
  \item[B.] BLACKHOLE
\end{itemize}
\column{0.48\textwidth}
\begin{itemize}\smallitem
  \item[C.] MyISAM
  \item[D.] ARCHIVE
\end{itemize}
\end{columns}
\end{frame}

% 31
\begin{frame}{Quick Quiz 10}
\begin{block}{Question 10}
If a table requires foreign key integrity, which engine is most suitable?
\end{block}
\vspace{0.2cm}
\begin{columns}[T,onlytextwidth]
\column{0.48\textwidth}
\begin{itemize}\smallitem
  \item[A.] InnoDB
  \item[B.] MEMORY
\end{itemize}
\column{0.48\textwidth}
\begin{itemize}\smallitem
  \item[C.] BLACKHOLE
  \item[D.] CSV
\end{itemize}
\end{columns}
\end{frame}

% 32
\begin{frame}{Quick Quiz 11}
\begin{block}{Question 11}
Which engine is well-suited for large, append-only, infrequently updated logs?
\end{block}
\vspace{0.2cm}
\begin{columns}[T,onlytextwidth]
\column{0.48\textwidth}
\begin{itemize}\smallitem
  \item[A.] ARCHIVE
  \item[B.] BLACKHOLE
\end{itemize}
\column{0.48\textwidth}
\begin{itemize}\smallitem
  \item[C.] FEDERATED
  \item[D.] MEMORY
\end{itemize}
\end{columns}
\end{frame}

% 26
\begin{frame}{Review Questions}
\begin{enumerate}\smallitem
  \item What purpose does a storage engine serve in MySQL?
  \item Why is \code{InnoDB} generally recommended for core business entities?
  \item Differentiate \code{InnoDB} and \code{MEMORY} regarding data persistence.
  \item Provide a realistic scenario where \code{ARCHIVE} is an optimal choice.
  \item Why should \code{BLACKHOLE} never be chosen for data requiring storage?
\end{enumerate}
\end{frame}

% 27
\begin{frame}{Practical Exercises 1--2}
\begin{block}{Exercise 1}
Write an SQL statement to list all storage engines available on your MySQL Server.
\end{block}
\begin{block}{Exercise 2}
Create a table named \code{students} with columns \code{student\_id}, \code{student\_name}, and \code{email}. The table must explicitly use \code{InnoDB}.
\end{block}
\begin{alertblock}{Requirement}
Use standard SQL syntax and explicitly specify the engine when creating the table.
\end{alertblock}
\end{frame}

% 28
\begin{frame}{Practical Exercises 3--4}
\begin{block}{Exercise 3}
Create a table named \code{temp\_report} with columns \code{report\_id}, \code{report\_name}, and \code{created\_at} using \code{MEMORY}. Explain its major operational limitation.
\end{block}
\begin{block}{Exercise 4}
An e-commerce system needs to store customers, orders, and payments. It requires transaction rollbacks on errors and strict foreign key integrity.
\end{block}
\begin{alertblock}{Requirement}
Select the most suitable storage engine and justify your technical rationale.
\end{alertblock}
\end{frame}

% 29
\begin{frame}{Lesson Summary}
\begin{itemize}\smallitem
  \item Storage engines determine how MySQL stores, formats, and retrieves data per table.
  \item \code{InnoDB} is the default engine and standard choice for modern business applications.
  \item \code{InnoDB} provides ACID transactions, foreign keys, row locking, and automatic crash recovery.
  \item \code{MEMORY} delivers high-speed in-RAM access but is volatile across server restarts.
  \item \code{ARCHIVE} is specialized for compressed, append-only historical log datasets.
  \item \code{CSV}, \code{BLACKHOLE}, \code{MERGE}, and \code{FEDERATED} address specialized technical needs.
\end{itemize}
\end{frame}

% 30
\begin{frame}{Key Terms}
\begin{columns}[T,onlytextwidth]
\column{0.33\textwidth}
\begin{itemize}\smallitem
  \item MySQL
  \item Storage engine
  \item InnoDB
  \item MyISAM
  \item MEMORY
\end{itemize}
\column{0.33\textwidth}
\begin{itemize}\smallitem
  \item ARCHIVE
  \item CSV
  \item BLACKHOLE
  \item MERGE
  \item FEDERATED
\end{itemize}
\column{0.33\textwidth}
\begin{itemize}\smallitem
  \item Transaction
  \item Foreign key
  \item \code{SHOW ENGINES}
  \item \code{CREATE TABLE}
  \item \code{ENGINE}
\end{itemize}
\end{columns}
\end{frame}

% 32
\begin{frame}{Quick Quiz Answer Key}
\begin{center}
\begin{tabular}{ccccccccccc}
\toprule
\textbf{Q1} & \textbf{Q2} & \textbf{Q3} & \textbf{Q4} & \textbf{Q5} & \textbf{Q6} & \textbf{Q7} & \textbf{Q8} & \textbf{Q9} & \textbf{Q10} & \textbf{Q11} \\
\midrule
B & B & C & B & B & C & A & A & B & A & A \\
\bottomrule
\end{tabular}
\end{center}
\vspace{0.35cm}
\begin{block}{Self-Check Guidance}
For any incorrect answer, review the corresponding slide to clarify the underlying concept: default engines, \code{ENGINE} syntax, durability, transactions, or foreign keys.
\end{block}
\end{frame}

% 33
\begin{frame}[fragile]{Exercise 1 Solution}
\begin{block}{Method 1}
\begin{lstlisting}
SHOW ENGINES;
\end{lstlisting}
\end{block}
\begin{block}{Method 2}
\begin{lstlisting}
SELECT
    engine,
    support
FROM information_schema.engines
ORDER BY engine;
\end{lstlisting}
\end{block}
\end{frame}

% 34
\begin{frame}[fragile]{Exercise 2 Solution}
\begin{lstlisting}
CREATE TABLE students (
    student_id INT PRIMARY KEY,
    student_name VARCHAR(100),
    email VARCHAR(100)
) ENGINE = InnoDB;
\end{lstlisting}
\begin{block}{Best Practice}
For production tables, explicitly declaring \code{ENGINE = InnoDB} avoids unexpected behavior if server default configurations change.
\end{block}
\end{frame}

% 33
\begin{frame}[fragile]{Exercise 3 Solution}
\begin{lstlisting}
CREATE TABLE temp_report (
    report_id INT PRIMARY KEY,
    report_name VARCHAR(100),
    created_at DATETIME
) ENGINE = MEMORY;
\end{lstlisting}
\begin{alertblock}{Limitation}
Data in \code{MEMORY} tables is volatile and will be lost whenever MySQL Server restarts or shuts down.
\end{alertblock}
\end{frame}

% 34
\begin{frame}{Exercise 4 Solution}
\begin{block}{Recommended Storage Engine}
Select \code{InnoDB}.
\end{block}
\begin{block}{Rationale}
\begin{itemize}\smallitem
  \item The workload requires ACID transactions with rollbacks upon execution errors.
  \item Orders must reference valid customers through foreign key constraints.
  \item Customer, order, and payment data require strict relational consistency.
  \item \code{InnoDB} provides crash recovery through redo logging for mission-critical business data.
\end{itemize}
\end{block}
\end{frame}

\end{document}
